% !TEX program = xelatex
% 编译：xelatex qcd_factorization.tex（运行三次以生成目录与交叉引用）
% 字体：fontset=fandol 随 TeX Live 发行，可在任何平台编译；
%       在 macOS 上可改为 fontset=macnew 以获得更好的中文字体效果。
\documentclass[UTF8,a4paper,11pt,oneside,openany,fontset=fandol]{ctexbook}

\usepackage[margin=2.5cm,headheight=14pt]{geometry}
\usepackage{amsmath,amssymb,amsthm,mathtools,bm}
% \mathbb 1：amssymb 的 msbm 字体没有数字 1，改用 dsfont 的 \mathds{1}；\mathbb{C}、\mathbb{R} 等不变
\usepackage{dsfont}
\let\amsmathbb\mathbb
\renewcommand{\mathbb}[1]{\ifx1#1\mathds{1}\else\amsmathbb{#1}\fi}
\usepackage{slashed}
\usepackage{booktabs,array,multirow,tabularx}
\usepackage{enumitem}
\usepackage{caption,subcaption}
\captionsetup{labelsep=quad}     % 图表标题 “图 5　……”：编号后空一字，不用西文冒号
\usepackage[section]{placeins}   % 浮动体不越过 \section 边界
\usepackage[dvipsnames]{xcolor}
\usepackage{tikz}
\usetikzlibrary{decorations.pathmorphing,decorations.markings,arrows.meta,calc,positioning,shapes.geometric,patterns,backgrounds,fit}
\usepackage{pgfplots}
\pgfplotsset{compat=1.17}
\usepackage[most]{tcolorbox}
\usepackage{etoolbox}            % \patchcmd、\appto（tcolorbox 亦载入，此处显式声明）
\usepackage[colorlinks=true,linkcolor=blue!50!black,citecolor=green!40!black,
            urlcolor=blue!60!black,bookmarksnumbered=true]{hyperref}
\usepackage[nameinlink]{cleveref}
\usepackage[nottoc]{tocbibind}   % 参考文献进入目录与 PDF 书签

% ---------------- 定理类环境 ----------------
% 定理、引理、命题、推论的正文：汉字用楷体，西文保持直立。plain 样式的 \itshape
% 会把对称性名称 “C 宇称/C 偶/C 奇” 中的正体 C 排成斜体，与电荷共轭矩阵 $C$
% 无法区分；上下间距与标题字体同 plain。
% 三种样式的标题都按中文习惯排：附注用全角括号（……），标题后不加西文句点而空一字；
% cdefinition、cremark 的字体与上下间距分别同 amsthm 的 definition、remark 样式。
\newtheoremstyle{cplain}{\topsep}{\topsep}{\kaishu}{0pt}{\bfseries}{}{1em}%
  {\thmname{#1}\thmnumber{ \textup{#2}}\thmnote{\normalfont（#3）}}
\newtheoremstyle{cdefinition}{\topsep}{\topsep}{\normalfont}{0pt}{\bfseries}{}{1em}%
  {\thmname{#1}\thmnumber{ \textup{#2}}\thmnote{\normalfont（#3）}}
\newtheoremstyle{cremark}{\glueexpr\topsep/2\relax}{\glueexpr\topsep/2\relax}%
  {\normalfont}{0pt}{\itshape}{}{1em}%
  {\thmname{#1}\thmnumber{ \textup{#2}}\thmnote{\normalfont（#3）}}
\theoremstyle{cplain}
\newtheorem{theorem}{定理}[chapter]
\newtheorem{lemma}[theorem]{引理}
\newtheorem{proposition}[theorem]{命题}
\newtheorem{corollary}[theorem]{推论}
\newtheorem{conjecture}[theorem]{猜想}
\theoremstyle{cdefinition}
\newtheorem{definition}[theorem]{定义}
\newtheorem{assumption}[theorem]{假设}
\newtheorem{example}[theorem]{例}
\theoremstyle{cremark}
\newtheorem{remark}[theorem]{注}
\renewcommand{\proofname}{证明}
% 证明的标题同样不加句点：amsthm 原定义中的 \@addpunct{.} 改为 0.5em，连同 \labelsep 共空一字
\makeatletter
\renewenvironment{proof}[1][\proofname]{\par\pushQED{\qed}\normalfont\topsep6\p@\@plus6\p@\relax
  \trivlist\item[\hskip\labelsep\itshape #1\hspace{0.5em}]\ignorespaces}{\popQED\endtrivlist\@endpefalse}
\makeatother

\tcolorboxenvironment{theorem}{enhanced,breakable,colback=blue!3!white,
  colframe=blue!45!black,boxrule=0.6pt,arc=2pt,left=5pt,right=5pt,top=3pt,bottom=3pt}

\tcolorboxenvironment{assumption}{enhanced,breakable,colback=orange!4!white,
  colframe=orange!65!black,boxrule=0.6pt,arc=2pt,left=5pt,right=5pt,top=3pt,bottom=3pt}
\tcolorboxenvironment{conjecture}{enhanced,breakable,colback=Plum!4!white,
  colframe=Plum!70!black,boxrule=0.6pt,arc=2pt,left=5pt,right=5pt,top=3pt,bottom=3pt}

\crefname{theorem}{定理}{定理}
\crefname{lemma}{引理}{引理}
\crefname{proposition}{命题}{命题}
\crefname{corollary}{推论}{推论}
\crefname{conjecture}{猜想}{猜想}
\crefname{assumption}{假设}{假设}
\crefname{definition}{定义}{定义}
\crefname{example}{例}{例}
\crefname{remark}{注}{注}
\crefname{equation}{式}{式}
\creflabelformat{equation}{#2(#1)#3}
\crefname{figure}{图}{图}
\crefname{table}{表}{表}
\crefformat{chapter}{#2第~#1~章#3}
\crefformat{section}{#2第~#1~节#3}
\crefformat{subsection}{#2第~#1~小节#3}
\crefformat{appendix}{#2附录~#1#3}
\crefformat{subappendix}{#2附录~#1#3}
\Crefformat{chapter}{#2第~#1~章#3}
\Crefformat{section}{#2第~#1~节#3}
\Crefformat{subsection}{#2第~#1~小节#3}
\Crefformat{appendix}{#2附录~#1#3}
\Crefformat{subappendix}{#2附录~#1#3}
% 多目标引用的中文连接词
\newcommand{\crefpairconjunction}{\CJKecglue 与\CJKecglue}
\newcommand{\crefmiddleconjunction}{、}
\newcommand{\creflastconjunction}{\CJKecglue 与\CJKecglue}
\newcommand{\crefpairgroupconjunction}{\CJKecglue 与\CJKecglue}
\newcommand{\crefmiddlegroupconjunction}{、}
\newcommand{\creflastgroupconjunction}{\CJKecglue 与\CJKecglue}
\newcommand{\crefrangeconjunction}{\CJKecglue 至\CJKecglue}

% \cref 后紧跟汉字时补中西文间距：xeCJK 不能越过超链接的边界自动插入 \CJKecglue
% （源文件中 \cref 与汉字之间不留空格），此处手工补上。
% 最后一组引用若为 section/subsection 类型（“第 n 节/小节”，以汉字结尾）则不补。
\ExplSyntaxOn
\NewCommandCopy \__fact_cref_orig:w \cref
\tl_new:N \l__fact_cref_type_tl
\seq_new:N \l__fact_cref_types_seq
% 标签 #1 的 cleveref 类型，取自 \r@#1@cref = {[类型][..][..]..}{..}；标签未定义时为空
\cs_new:Npn \__fact_cref_type:n #1
  {
    \cs_if_exist:cT { r@ #1 @cref }
      { \exp_last_unbraced:Nv \__fact_cref_type_aux:w { r@ #1 @cref } \q_stop }
  }
\cs_new:Npn \__fact_cref_type_aux:w #1#2 \q_stop
  { \__fact_cref_type_auxii:w #1 [] \q_stop }
\cs_new:Npn \__fact_cref_type_auxii:w #1 [#2] #3 \q_stop
  { \tl_if_empty:nT {#1} {#2} }
% 输出中最后一组引用的类型（cleveref 默认按各类型首次出现的次序分组输出）
\cs_new_protected:Npn \__fact_cref_last_type:n #1
  {
    \seq_clear:N \l__fact_cref_types_seq
    \clist_map_inline:nn {#1}
      {
        \tl_set:Nx \l__fact_cref_type_tl { \__fact_cref_type:n {##1} }
        \seq_if_in:NVF \l__fact_cref_types_seq \l__fact_cref_type_tl
          { \seq_put_right:NV \l__fact_cref_types_seq \l__fact_cref_type_tl }
      }
    \seq_get_right:NNF \l__fact_cref_types_seq \l__fact_cref_type_tl
      { \tl_clear:N \l__fact_cref_type_tl }
  }
% 下一记号若为 CJK 统一汉字（U+3400--U+9FFF），在其前插入 \CJKecglue
\cs_new_protected:Npn \__fact_cref_ecglue:N #1
  {
    \bool_lazy_all:nTF
      {
        { \token_if_letter_p:N #1 }
        { ! \token_if_cs_p:N #1 }
        { \int_compare_p:nNn { `#1 } > { "33FF } }
        { \int_compare_p:nNn { `#1 } < { "A000 } }
      }
      { \CJKecglue #1 }
      { #1 }
  }
\RenewDocumentCommand \cref { s m }
  {
    \IfBooleanTF {#1}
      { \__fact_cref_orig:w * {#2} }
      { \__fact_cref_orig:w {#2} }
    \__fact_cref_last_type:n {#2}
    \str_case:VnF \l__fact_cref_type_tl
      { { chapter } { } { section } { } { subsection } { } }
      { \peek_N_type:TF { \__fact_cref_ecglue:N } { } }
  }
\ExplSyntaxOff

\numberwithin{equation}{chapter}

% 标题中的数学式随标题加粗（含目录中节一级条目）
\ctexset{chapter/number=\arabic{chapter}}   % 章号用阿拉伯数字：“第 5 章”、式 (5.3)
\ctexset{chapter/format+=\boldmath,section/format+=\boldmath,subsection/format+=\boldmath,paragraph/format+=\boldmath}
\makeatletter
\patchcmd{\l@section}{\bfseries}{\bfseries\boldmath}{}%
  {\@latex@warning@no@line{\string\l@section\space not patched: math in TOC section entries stays non-bold}}
\makeatother

% 目录行距略紧并单独成页（正文从新页开始）
\let\factorigtableofcontents\tableofcontents
\renewcommand{\tableofcontents}{%
  \begingroup\linespread{1.1}\selectfont\factorigtableofcontents\endgroup\clearpage}

% 参考文献条目间距略紧。\parsep 须在列表参数中设定（\list 随即把它赋给 \parskip），
% 故写入 thebibliography 在列表参数中执行的 \@openbib@code，而不追加在 \thebibliography 之后
\makeatletter
\appto\@openbib@code{\setlength{\itemsep}{1pt plus 1pt}\setlength{\parsep}{0pt}}
\makeatother

% 页眉：节标题不转为大写（避免 “DIRAC”“FURRY” 之类）；页眉一律用直立体（含目录、
% 参考文献页）：headings 样式原以 \slshape 排页眉，而 FandolSong 没有倾斜体，汉字会
% 退成楷体，西文则倾斜
\makeatletter
\renewcommand{\chaptermark}[1]{\markright{\ifnum\c@secnumdepth>\m@ne\CTEXthechapter\quad\fi #1}}
\renewcommand{\sectionmark}[1]{}
\def\@oddhead{{\upshape\rightmark}\hfil\thepage}
\makeatother

% ---------------- 记号宏 ----------------
\DeclareMathOperator{\Tr}{Tr}
\DeclareMathOperator{\Real}{Re}
\DeclareMathOperator{\Imag}{Im}
\DeclareMathOperator{\Disc}{Disc}
\newcommand{\TrD}{\operatorname{Tr}_{\mathrm{D}}}   % Dirac 迹
\newcommand{\Trc}{\operatorname{Tr}_{\mathrm{c}}}   % 色迹
\newcommand{\sla}[1]{\slashed{#1}}                  % Feynman slash
\newcommand{\dd}{\mathrm{d}}                         % 微分符号
\newcommand{\tp}{\mathsf{T}}                         % 转置
\newcommand{\ii}{\mathrm{i}}                        % 虚数单位（正体）
\newcommand{\ee}{\mathrm{e}}                        % 自然对数底（正体）
\newcommand{\MSbar}{\overline{\mathrm{MS}}}         % MS-bar 方案
\newcommand{\order}[1]{\mathcal{O}\!\left(#1\right)}% 量级
\newcommand{\Lqcd}{\Lambda_{\mathrm{QCD}}}
\newcommand{\as}{\alpha_s}
\newcommand{\pT}{p_{\mathrm{T}}}
\newcommand{\kT}{k_{\mathrm{T}}}
\newcommand{\bT}{b_{\mathrm{T}}}
\newcommand{\qT}{q_{\mathrm{T}}}

% ---------------- 严格性标签 ----------------
% [R] 严格  [A] 全阶论证  [P] 部分  [C] 猜想  [V] 破坏
\newtcbox{\facttagbox}[1][gray]{on line,arc=2pt,boxsep=0pt,left=2.2pt,right=2.2pt,top=1.2pt,bottom=1.2pt,
  boxrule=0.45pt,colback=#1!10!white,colframe=#1!70!black,fontupper=\footnotesize\sffamily\upshape}
\newcommand{\tagR}{\facttagbox[OliveGreen]{R\,严格}}
\newcommand{\tagA}{\facttagbox[RoyalBlue]{A\,全阶论证}}
\newcommand{\tagP}{\facttagbox[BurntOrange]{P\,部分}}
\newcommand{\tagC}{\facttagbox[Plum]{C\,猜想}}
\newcommand{\tagV}{\facttagbox[BrickRed]{V\,破坏}}
% 依赖记录（仿 Lean blueprint 的 \uses）：在定义/假设/引理/定理等环境内部写 \uses{lab1,lab2}，
% 正文中以小号字列出所依赖的条目；附录中的“形式化蓝图”由脚本据此自动生成依赖表。
\newcommand{\uses}[1]{\par\noindent{\footnotesize\color{gray!45!black}依赖：\cref{#1}}\par}
% 只给出证明要点（对 [A]/[P] 结论）时使用
\newenvironment{proofsketch}[1][证明概要]{\begin{proof}[#1]}{\end{proof}}

% ---------------- TikZ 费曼图样式（纯 TikZ，不依赖 tikz-feynman） ----------------
%%% BEGIN STYLES
% =====================================================================
%  Plain-TikZ Feynman-diagram styles (no tikz-feynman).
%  Load in the preamble, after \usepackage{tikz}:
%    \usetikzlibrary{decorations.pathmorphing,decorations.markings,arrows.meta,calc,positioning}
%  Contents of this block (all low-level names start with \fd):
%   (i)  helper registers/macros and three decorations
%          "fd gluon"      coil for open paths,
%          "fd gluon loop" coil for closed paths (full circles),
%          "fd photon"     sine wave;
%        they rescale the nominal segment length so that an integer
%        number of coils / half-waves fits the path exactly (no straight
%        left-overs, evenly spaced coils, also on arcs);
%   (ii) a single \tikzset with the user-level styles:
%          fd             picture style (line width, caps, label font)
%          quark[=pos]    solid line, centred Stealth arrow at pos (default 0.5)
%          gluon          coiled line;   gluon loop[=m]: closed coil, number
%                         of coils a multiple of m (use for full circles)
%          photon         wavy line
%          ghost[=pos]    dotted line with centred arrow
%          mom={..}, mom'={..}   momentum arrow parallel to a path, on its
%                         left / right:  \path[mom={$p$}] (A) -- (B);
%                         keys: mom pos, mom length, mom sep
%          vertex         vertex dot (node style):  \node[vertex] at (A) {};
%          blob           grey filled region:  \draw[blob] (c) circle[radius=5mm];
%          bound state    double line (e.g. a bound state)
% =====================================================================
\newdimen\fdAA \newdimen\fdTW \newdimen\fdHW
\def\fdloopmult{1}
\def\fdpt#1#2#3{\pgfqpoint{\dimexpr#1\fdAA+#3\fdTW\relax}{#2\pgfdecorationsegmentamplitude}}
\def\fdcoilfull{%
  \pgfpathcurveto{\fdpt{0}{0.555}{1}}{\fdpt{0.445}{1}{2}}{\fdpt{1}{1}{3}}%
  \pgfpathcurveto{\fdpt{1.555}{1}{4}}{\fdpt{2}{0.555}{5}}{\fdpt{2}{0}{6}}%
  \pgfpathcurveto{\fdpt{2}{-0.555}{7}}{\fdpt{1.555}{-1}{8}}{\fdpt{1}{-1}{9}}%
  \pgfpathcurveto{\fdpt{0.445}{-1}{10}}{\fdpt{0}{-0.555}{11}}{\fdpt{0}{0}{12}}}
\def\fdcoilhalf{%
  \pgfpathcurveto{\fdpt{0}{0.555}{1}}{\fdpt{0.445}{1}{2}}{\fdpt{1}{1}{3}}%
  \pgfpathcurveto{\fdpt{1.555}{1}{4}}{\fdpt{2}{0.555}{5}}{\fdpt{2}{0}{6}}}
\def\fdcoilopeninit{%
  \fdAA=\pgfdecorationsegmentaspect\pgfdecorationsegmentamplitude\relax
  \pgfmathsetmacro\fdN{max(1,round((\pgfdecoratedpathlength-2*\fdAA)/\pgfdecorationsegmentlength-0.5))}%
  \pgfmathsetlength\pgfdecorationsegmentlength{(\pgfdecoratedpathlength-2*\fdAA-0.3pt)/(\fdN+0.5)}%
  \fdTW=0.0833333\pgfdecorationsegmentlength\relax}%
\def\fdcoilclosedinit{%
  \fdAA=\pgfdecorationsegmentaspect\pgfdecorationsegmentamplitude\relax
  \pgfmathsetmacro\fdN{\fdloopmult*max(1,round(\pgfdecoratedpathlength/(\fdloopmult*\pgfdecorationsegmentlength)))}%
  \pgfmathsetlength\pgfdecorationsegmentlength{(\pgfdecoratedpathlength-0.1pt)/\fdN}%
  \fdTW=0.0833333\pgfdecorationsegmentlength\relax}%
\def\fdphotoninit{%
  \pgfmathsetmacro\fdN{max(1,round(2*\pgfdecoratedpathlength/\pgfdecorationsegmentlength))}%
  \pgfmathsetlength\fdHW{(\pgfdecoratedpathlength-0.1pt)/\fdN}}%
\pgfdeclaredecoration{fd gluon}{init}{%
  \state{init}[width=+0pt,next state=coil,persistent precomputation=\fdcoilopeninit]{}%
  \state{coil}[switch if less than=1.5\pgfdecorationsegmentlength+2\fdAA to last,
               width=+\pgfdecorationsegmentlength]{\fdcoilfull}%
  \state{last}[width=+\dimexpr.5\pgfdecorationsegmentlength+2\fdAA\relax,next state=final]{\fdcoilhalf}%
  \state{final}{\pgfpathlineto{\pgfpointdecoratedpathlast}}}
\pgfdeclaredecoration{fd gluon loop}{init}{%
  \state{init}[width=+0pt,next state=coil,persistent precomputation=\fdcoilclosedinit]{}%
  \state{coil}[switch if less than=+\pgfdecorationsegmentlength to final,
               width=+\pgfdecorationsegmentlength]{\fdcoilfull}%
  \state{final}{\pgfpathlineto{\pgfpointdecoratedpathlast}}}
\pgfdeclaredecoration{fd photon}{init}{%
  \state{init}[width=+0pt,next state=up,persistent precomputation=\fdphotoninit]{}%
  \state{up}[switch if less than=+\fdHW to final,width=+\fdHW,next state=down]{%
    \pgfpathsine{\pgfqpoint{.5\fdHW}{\pgfdecorationsegmentamplitude}}%
    \pgfpathcosine{\pgfqpoint{.5\fdHW}{-\pgfdecorationsegmentamplitude}}}%
  \state{down}[switch if less than=+\fdHW to final,width=+\fdHW,next state=up]{%
    \pgfpathsine{\pgfqpoint{.5\fdHW}{-\pgfdecorationsegmentamplitude}}%
    \pgfpathcosine{\pgfqpoint{.5\fdHW}{\pgfdecorationsegmentamplitude}}}%
  \state{final}{\pgfpathlineto{\pgfpointdecoratedpathlast}}}
\tikzset{
  /pgf/decoration/coil multiple/.store in=\fdloopmult,
  fd/.style={line width=0.8pt,line cap=round,line join=round,font=\small},
  fd arrow/.style={postaction={decorate,solid,line width=0.8pt,decoration={markings,
      mark=at position #1 with {\arrow[xshift=2.9pt]{Stealth[length=5.5pt,width=4.5pt]}}}}},
  quark/.style={draw,solid,line width=0.8pt,fd arrow=#1},
  quark/.default=0.5,
  plain/.style={draw,solid,line width=0.8pt},
  gluon/.style={draw,solid,line width=0.75pt,decorate,
      decoration={fd gluon,amplitude=3.2pt,segment length=4.6pt,aspect=0.7}},
  gluon loop/.style={draw,solid,line width=0.75pt,decorate,
      decoration={fd gluon loop,amplitude=3.2pt,segment length=4.6pt,aspect=0.7,coil multiple=#1}},
  gluon loop/.default=1,
  photon/.style={draw,solid,line width=0.8pt,decorate,
      decoration={fd photon,amplitude=2pt,segment length=6.6pt}},
  ghost/.style={draw,line width=1.3pt,line cap=round,dash pattern=on 0pt off 2.6pt,fd arrow=#1},
  ghost/.default=0.5,
  mom pos/.store in=\fdmompos,mom pos=0.5,
  mom length/.store in=\fdmomlen,mom length=0.55cm,
  mom sep/.store in=\fdmomsep,mom sep=0.25cm,
  mom arrow/.style={draw,solid,line width=0.5pt,-{Stealth[length=3.8pt,width=3.2pt]}},
  mom/.style={decorate,decoration={markings,mark=at position \fdmompos with {%
      \draw[mom arrow] (-0.5*\fdmomlen,\fdmomsep) -- node[auto,inner sep=1.5pt]{#1} (0.5*\fdmomlen,\fdmomsep);}}},
  mom'/.style={decorate,decoration={markings,mark=at position \fdmompos with {%
      \draw[mom arrow] (-0.5*\fdmomlen,-\fdmomsep) -- node[auto,swap,inner sep=1.5pt]{#1} (0.5*\fdmomlen,-\fdmomsep);}}},
  vertex/.style={circle,fill,draw=none,inner sep=0pt,minimum size=3.8pt},
  blob/.style={draw,solid,fill=gray!30,line width=0.8pt},
  bound state/.style={draw,solid,double,double distance=1.6pt,line width=0.6pt},
}
%%% END STYLES

%%% BEGIN STYLES-EXT
% =====================================================================
%  Extension of the plain-TikZ Feynman-diagram library (fd_styles.tex)
%  for the factorization monograph.  Requires fd_styles.tex loaded first.
%  Extra libraries: shapes.geometric, patterns, backgrounds, fit.
%  New user-level styles:
%    wilson[=pos]   Wilson (eikonal) line: double line with centred arrow
%    wilson plain   Wilson line without arrow
%    cut            final-state cut: dashed red vertical line
%    hard blob      hard subgraph (grey);   jet blob (A side, blue);
%    jetB blob      collinear subgraph of hadron B (orange);
%    soft blob      soft subgraph (green);  hadron blob (target, light grey)
%    glauber        Glauber gluon (dark-red coil)
%    opvertex       operator insertion vertex (circled cross), node style
%    scalar         dashed line (e.g. scalar/longitudinal polarization)
%    longit         gluon carrying k^mu (scalar) polarization: blue coil;
%                   optionally mark it with \longitmark{(A)}{(B)} (dot at midpoint)
%    lcaxis         thin arrow for light-cone / spacetime axes
%    region         dashed rounded rectangle marking a momentum region
% Colours are kept light so that the figures print well in greyscale.
% =====================================================================
\tikzset{
  wilson/.style={draw,solid,double,double distance=1.6pt,line width=0.55pt,
      postaction={decorate,solid,decoration={markings,mark=at position #1 with
      {\arrow[xshift=3.2pt,line width=0.6pt]{Stealth[length=6.4pt,width=6pt]}}}}},
  wilson/.default=0.5,
  wilson plain/.style={draw,solid,double,double distance=1.6pt,line width=0.55pt},
  cut/.style={draw,dashed,red!70!black,line width=0.9pt,dash pattern=on 4pt off 2.5pt},
  hard blob/.style={draw,solid,fill=gray!35,line width=0.8pt},
  jet blob/.style={draw,solid,fill=blue!12,line width=0.8pt},
  jetB blob/.style={draw,solid,fill=orange!18,line width=0.8pt},
  soft blob/.style={draw,solid,fill=green!15,line width=0.8pt},
  hadron blob/.style={draw,solid,fill=gray!15,line width=0.8pt},
  glauber/.style={draw,solid,line width=0.75pt,color=red!60!black,decorate,
      decoration={fd gluon,amplitude=3.2pt,segment length=4.6pt,aspect=0.7}},
  opvertex/.style={circle,draw,solid,fill=white,line width=0.8pt,inner sep=0pt,minimum size=7pt,
      path picture={\draw[line width=0.6pt] (path picture bounding box.north west)--(path picture bounding box.south east)
                    (path picture bounding box.south west)--(path picture bounding box.north east);}},
  scalar/.style={draw,solid,line width=0.8pt,dash pattern=on 3pt off 2pt},
  longit/.style={draw,solid,line width=0.75pt,color=blue!60!black,decorate,
      decoration={fd gluon,amplitude=3.2pt,segment length=4.6pt,aspect=0.7}},
  lcaxis/.style={draw,solid,line width=0.5pt,-{Stealth[length=4pt,width=3.4pt]}},
  region/.style={draw,dashed,rounded corners=3pt,line width=0.6pt,gray!70!black},
}
\newcommand{\longitmark}[2]{\fill[blue!60!black] ($#1!0.5!#2$) circle[radius=1.7pt];}
%%% END STYLES-EXT

